\section{实验复现细节}\label{app:experiments}

配套脚本 \texttt{verification\_experiments.py} 执行正文中的统计量计算和递推过程。运行环境需要 Python 和 NumPy；受控实验不依赖外部市场数据或行为数据。随机种子、全部参数、原始轨迹数量、估计器汇总结果和逐期记录均保存为 JSON 与 CSV 文件。表格条目由这些已保存的输出生成，并采用与文中汇总结果相同的数值精度。估计器与学习过程的随机种子为 20260911，实际执行收益流的随机种子为 $88177$。逐层诊断在每个层级使用 $1600$ 次独立重复，批量诊断使用 $8192$ 次带截断的重复模拟。各批量结果只在所报告的不同分组规模之间复用同一个已保存样本池，回放中的学习样本则重新生成。带有限层级上限的数值模拟只表示相应的截断构造，不表示未截断模拟。

\subsection{多项式权重诊断中的解析偏差}
记 $b=1-2\eps$、$D_w=(1-\eps)^3-\eps^3$，则
\[
 b_{w,0}=3\eps^2b/D_w,\qquad b_{w,1}=6\eps b^2/D_w,
 \qquad b_{w,2}=3b^3/D_w.
\]
在纯收益诊断中取 $\theta=0$，有 $S(v_+(ct))=1-t$。直接对 $1+\log u$ 从 $t$ 积分到 $1$，得到
$\E[(1+\log U)\1\{U>t\}]=-t\log t$。外层样本与内层样本相互独立，并且 $\E G=0$，所以 $\E T_k^{(n)}$ 中只含 $(w^\eps)'(\widehat S)\widehat S$ 的项被消去。Bernoulli 经验均值满足 $\E\widehat S=S$ 和 $\E\widehat S^2=S^2+S(1-S)/n$，因此
\[
 \E (w^\eps)'(\widehat S)=(w^\eps)'(S)+b_{w,2}S(1-S)/n.
\]
将上式乘以 $\E[G\1\{U>t\}]$，作变量代换 $z=v_+(ct)$，再积分，即可证明式~\eqref{eq:diagnostic-bias}。所有积分均有限，因为固定的光滑价值函数使端点处的对数因子可积。若层级上限为 $L_{\max}$，则在样本量记号下，截断构造以 $T_k^{(N_02^{L_{\max}})}$ 为目标，由此得到文中报告的剩余偏差。这些公式为实际抽样统计量提供独立的总体比较基准，不会以解析量替代样本数据。

\subsection{有限样本统计量的精确积分}
对收益侧或损失侧中的一侧，先将零点、全部内层效用观测值和外层效用观测值合并，排序后去除重复值。在任意两个相邻断点之间，经验生存概率和外层示性函数都保持常数。因此，可用区间长度乘以 $(w^\eps)'(\widehat S)(I_0-\widehat S)$，再对全部区间求和。超过最大断点后，两者都为零；由于 $f(0,0)=0$，该区间上的被积函数也为零。对另一侧重复相同步骤，再取收益侧与损失侧之差，并乘以共享得分。两个半样本统计量使用相同的外层效用值和得分，但各自使用对应内层半样本产生的断点。

\subsection{独立数值梯度检验}
单期回放包含一只风险资产，其收益率为 $+r_*$ 或 $-r_*$，并且 $U\sim\operatorname{Beta}(e^\theta,1)$。固定情境 $(X,r)$ 后，正相对结果阈值 $y$ 等价于终端财富超过 $r+y/(1-\eta_+)$；损失阈值等价于终端财富低于 $r-y/(1-\eta_-)$。对 $R\ne0$，相应的动作阈值为 $(\text{wealth threshold}/X-1)/R$，并截断到 $[0,1]$。该阈值处的分布函数为 $t^{e^\theta}$，在区间内部对参数的导数为 $e^\theta t^{e^\theta}\log t$。合并两个收益率分支后，可直接得到 $S_\pm$ 和 $\nabla S_\pm$。数值计算分别对 $w^\eps(S_\pm)$ 和 $(w^\eps)'(S_\pm)\nabla S_\pm$ 关于效用阈值积分，并在端点断点处分割积分区间。另行对独立积分得到的目标使用中心差分，以交叉检验梯度。将积分阶数加倍得到正文报告的数值稳定性诊断；该比较不构成严格的积分误差证明。

\subsection{回放步长的适用区域}
对已执行的受控回放，$h=1$、$c_{\rm tc}=0$、$|R|\le0.04$、$K=36$，且 $\Theta=[-2,2]$。均值动作和抽样动作都位于 $[0,1]$，所以回放中所有期初财富以及单步假想财富均由 $X_{\max}=1.04^{K+1}$ 控制。参考点是此前财富与初始参考点的凸组合，因而同样满足这一上界。由非负性可得 $|Y|\le X_{\max}$。由于该验证中的两个指数和两个平滑尺度分别相同，取 $B_v=\lambda v_+(X_{\max})$。标量 Beta 得分为 $G=1+e^\theta\log U$；随机变量 $-e^\theta\log U$ 服从单位指数分布，因此 $\E G^2=1$ 且 $\E|\partial_\theta G|=1$。于是，由似然 Hessian 的计算可知，$L=2B_v(C_2+2C_1)$ 是一个有效的保守上界；实现中取 $a_0=1$、$\gamma=1/(2L)$，满足 $\gamma L\le a_0$。

\subsection{数值精度}
统计实验采用有限层级上限和浮点运算。报告的标准误用于量化 Monte Carlo 变动，积分阶数比较用于评估数值稳定性。解析结果对应满足既定假设的数学抽样构造。发布的源文件记录了随机种子、精度设置、重复模拟输出和逐期记录，从而可以直接比较数值量与解析量。
